\documentclass[11pt,a4paper]{article}

\usepackage[margin=1in]{geometry}
\usepackage[T1]{fontenc}
\usepackage{lmodern}
\usepackage{booktabs}
\usepackage{array}
\usepackage{amsmath}
\usepackage{xcolor}
\usepackage{listings}
\usepackage{enumitem}
\usepackage[hidelinks]{hyperref}
\usepackage{titlesec}

\titleformat{\section}{\normalfont\large\bfseries}{\thesection}{0.6em}{}
\titleformat{\subsection}{\normalfont\normalsize\bfseries}{\thesubsection}{0.6em}{}
\setlist[itemize]{leftmargin=1.4em,itemsep=2pt,topsep=3pt}
\setlist[enumerate]{leftmargin=1.6em,itemsep=2pt,topsep=3pt}

\definecolor{codebg}{gray}{0.96}
\lstset{
  basicstyle=\ttfamily\footnotesize,
  backgroundcolor=\color{codebg},
  frame=single,
  framerule=0pt,
  breaklines=true,
  columns=fullflexible,
  keepspaces=true,
  xleftmargin=0.5em,
  showstringspaces=false,
  aboveskip=0.7em,
  belowskip=0.7em
}

\newcommand{\code}[1]{\texttt{\small #1}}

\title{\vspace{-2em}\textbf{LoMix: Core Contributions, Replication Findings,}\\
\textbf{and Where It Goes Next}\\[0.4em]
\large Deliverable 2 --- Technical Report}
\author{Md. Irtiza Hossain \\
\small M.S. Computer Science \& Engineering, BRAC University \\
\small \href{mailto:mohammad.irtiza.hossain@gmail.com}{mohammad.irtiza.hossain@gmail.com}}
\date{6 September 2026}

\begin{document}
\maketitle

\noindent\begin{tabular}{@{}ll@{}}
\textbf{Paper} & Md Mostafijur Rahman, Radu Marculescu. \emph{LoMix: Learnable Weighted} \\
               & \emph{Multi-Scale Logits Mixing for Medical Image Segmentation.} NeurIPS 2025. \\
\textbf{Code}  & \code{github.com/SLDGroup/LoMix} @ \code{464a945} \\
\textbf{Data}  & Synapse Multi-organ (BTCV), 9 classes, $224\times224$ \\
\textbf{Hardware} & GTX 960, 4\,GB, capability 5.2 (Maxwell, no tensor cores) \\
\textbf{Harness} & \url{https://github.com/irtiza1999/LoMix-Reproducibility-Study} \\
\end{tabular}

\vspace{0.6em}
\noindent{\small Full setup notes, run logs, per-arm curves and reproduction
commands are in the accompanying execution summary (Deliverable 1). Every
number below is regenerable from the logs in the harness repository above.}

\section{What the paper does}

A multi-scale decoder emits one logit map per stage. Classical deep supervision
applies a loss to each and sums them with fixed, usually uniform weights. LoMix
changes two things.

\textbf{Mixing.} It supervises not only the individual maps but \emph{combinations}
of them. For four decoder outputs it enumerates all 11 subsets of size $\geq 2$
and fuses each subset under four operations --- \code{add}, \code{mul},
\code{wf} (a small learned fusion module) and \code{concat} --- producing 44
synthetic logit maps, each receiving its own CE + Dice loss.

\textbf{Learned weights.} Every individual map and every mixed map contributes
through a learned scalar passed through \code{softplus}, optimised jointly with
the network by the same AdamW step.

\begin{lstlisting}
  encoder -- decoder --+-- p4 --+
                       |-- p3 --|   subsets (11)  x  ops (add, mul, wf, concat)
                       |-- p2 --|        |
                       +-- p1 --+        v
                          |        44 mixed maps -- CE+Dice -- x softplus(w_op,subset)
                          |                                          |
                          +-------- CE+Dice -- x softplus(w_i) ------+--> total loss

  inference:  encoder -- decoder -- p1        (mixing lives only in the loss)
\end{lstlisting}

The property that makes this attractive in practice: \textbf{all of it lives in
the loss}. The mixing weights are parameters of
\code{CombinatorialMutationsLossModule}, not of \code{EMCADNet}. At test time
the network is unchanged --- same graph, same FLOPs, same latency. Training cost
rises; inference cost does not.

\section{Why I chose this paper}

It asks a question my own work keeps running into, in a different domain.
\textbf{CLEAR-MoE} (under review) extracts sparse experts from a \emph{frozen}
Vision Transformer and selects which layers to convert using a
calibration-driven criterion --- a decision about which internal signals to
trust. \textbf{Predictive Self-Conditioning} (under review, WACV 2027) feeds a
LiDAR detector's own predictions back as conditioning and is evaluated on
\emph{calibration}, not accuracy alone, because a self-conditioning loop that
trusts a wrong prediction compounds its own error. LoMix asks the same question
at the supervision level: \emph{which of my own multi-scale outputs deserve
weight?} It answers with a learned scalar.

\section{Replication design}

The repository exposes \code{-{}-supervision \{last\_layer, deep\_supervision,
mutation, lomix\}} over one fixed network, which is a ready-made ablation
ladder: the network, the data and the schedule are identical and only the loss
differs.

\begin{center}
\begin{tabular}{@{}ll@{}}
\toprule
\textbf{Arm} & \textbf{Isolates} \\
\midrule
\code{last\_layer} & no deep supervision at all \\
\code{deep\_supervision} & per-scale losses, uniform weights \\
\code{mutation} & + logit mixing, \textbf{unweighted} \\
\code{lomix} & + \textbf{learnable} weights \\
\bottomrule
\end{tabular}
\end{center}

\code{mutation} $\rightarrow$ \code{lomix} is the load-bearing comparison: it
isolates the learned weighting from the mixing itself. If the gain came only
from having more loss terms, that pair would show it.

Constraints forced a \textbf{reduced-budget reproduction}: 20 of the
repository's default 300 epochs (the paper's main text states no epoch budget), per-step batch 2 with gradient accumulation 3 to restore the paper's
effective batch of 6, fp32, validation every 2 epochs. Configuration was chosen
by measurement, not guesswork --- a memory/latency sweep showed that AMP is
2.8--3.2$\times$ \emph{slower} on a card without tensor cores, and that the
binding constraint is the Windows 2-second GPU watchdog rather than VRAM.

\section{Findings}

\subsection{The zero-inference-overhead claim holds}

\begin{center}
\begin{tabular}{@{}lrrrrr@{}}
\toprule
\textbf{Encoder} & \textbf{Params (M)} & \textbf{Dec.\ (M)} & \textbf{GMACs} &
\textbf{Latency} & \textbf{Peak VRAM} \\
\midrule
pvt\_v2\_b0 & 3.921 & 0.507 & 0.657 & 17.7\,ms & 34\,MiB \\
pvt\_v2\_b2 & 26.773 & 1.914 & 4.324 & 36.8\,ms & 147\,MiB \\
\bottomrule
\end{tabular}
\end{center}

All four arms share a byte-identical inference graph, so one measurement settles
it for all of them. The decoder's 1.914\,M parameters match EMCAD's published
decoder cost, confirming the model is built as intended.

\subsection{Training cost, which the paper does not report}

Same network, same batch, only \code{-{}-supervision} differs:

\begin{center}
\begin{tabular}{@{}lrrr@{}}
\toprule
\textbf{Arm} & \textbf{Peak train VRAM} & \textbf{Step time} &
\textbf{vs \code{last\_layer}} \\
\midrule
\code{last\_layer} & 1008\,MiB & 649\,ms & $1.00\times$ \\
\code{deep\_supervision} & 1054\,MiB & 708\,ms & $1.09\times$ \\
\code{mutation} & 1268\,MiB & 938\,ms & $1.45\times$ \\
\code{lomix} & 2290\,MiB & 1649\,ms & $\mathbf{2.54\times}$ \\
\bottomrule
\end{tabular}
\end{center}

LoMix costs $2.54\times$ step time and $2.27\times$ training memory, and two
thirds of that gap comes from the learnable-weight arm alone. On constrained
hardware this decides whether the method is runnable at all --- the same class
of concern the lab's efficiency work targets, applied to training rather than
inference.

\subsection{Accuracy: the ordering did not reproduce at 20 epochs}

\begin{center}
\begin{tabular}{@{}lrl@{}}
\toprule
\textbf{Arm} & \textbf{Best mean Dice} & \textbf{Status} \\
\midrule
\code{last\_layer} & 0.6764 & completed \\
\code{deep\_supervision} & \textbf{0.7979} & completed \\
\code{mutation} & 0.7788 & completed \\
\code{lomix} & 0.7225 (pre-divergence) & \textbf{diverged (NaN) at epoch 13} \\
\bottomrule
\end{tabular}
\end{center}

The paper's Table~1 reports these same four arms on this same backbone
(PVT-EMCAD-B2, Synapse 8-organ), averaged over at least three runs on an RTX
A6000, which makes the comparison direct --- though not like-for-like:

\begin{center}
\begin{tabular}{@{}lrrrcr@{}}
\toprule
& \multicolumn{3}{c}{\textbf{Paper (Table 1)}} & & \textbf{This work} \\
\cmidrule(lr){2-4}\cmidrule(l){6-6}
\textbf{Arm} & \textbf{DICE} & \textbf{HD95} & \textbf{mIoU} & & \textbf{DICE} \\
\midrule
LL (\code{last\_layer}) & 80.9 & 22.9 & 71.2 & & 67.6 \\
DS (\code{deep\_supervision}) & 82.9 & 19.7 & 73.8 & & \textbf{79.8} \\
\code{mutation} & 83.6 & 15.7 & 74.7 & & 77.9 \\
\code{lomix} & \textbf{85.1} & \textbf{14.9} & \textbf{76.4} & & 72.3$^\dagger$ \\
\bottomrule
\end{tabular}
\end{center}
{\small Percentages. $^\dagger$best before divergence at epoch 13, not a
completed run.}

The published ordering is monotone; mine is not. My absolute values are
uniformly lower, and my best arm (79.8) sits just below the paper's
\emph{weakest} arm (80.9) --- what 20 epochs against a converged schedule looks
like, rather than a broken setup. The margins at stake are also small: the
paper's DS $\rightarrow$ LoMix gap is 2.2 DICE, established over at least three
runs with a two-sided Wilcoxon signed-rank test. One seed at 20 epochs cannot
resolve an effect that size.

Two further observations. \code{last\_layer} peaks at its \emph{first}
validation and decays thereafter, while the supervised arms climb steadily --- consistent with
\S4.5. And the learned mixing weights had barely differentiated: from a common
initialisation of $\mathrm{softplus}(0) = 0.693$, the four per-scale weights
ended at 0.379, 0.386, 0.392, 0.400. The mechanism supposed to create LoMix's
advantage had not yet engaged at this budget, which is the most likely
explanation for the ordering and is not separable from it with one seed at 20
epochs.

\subsection{The \code{lomix} arm diverged, and the cause is in the loss}

At iteration 14900, epoch 13, all loss terms became NaN simultaneously and every
later validation returned 0.000. Not a hardware event: fp32 throughout, no OOM,
no watchdog reset, and the saved weights are finite. Two defects in the
combinatorial loss explain it, both verified empirically rather than read off
the source (\code{tools/diagnose\_loss\_ops.py}).

\textbf{\code{mul} multiplies raw logits}, so a $k$-map combination is a
$k$-fold product growing like $|\text{logit}|^k$. Measured on the trained
checkpoint's own logits ($\max |p| = 6.16$):

\begin{center}
\begin{tabular}{@{}lrrr@{}}
\toprule
\textbf{Operation} & $k=2$ & $k=3$ & $k=4$ \\
\midrule
\code{add} & 12.0 & 17.9 & 23.6 \\
\textbf{\code{mul}} & \textbf{36.2} & \textbf{211.0} & \textbf{1214.0} \\
\code{wf} & 6.1 & 6.1 & 5.9 \\
\code{concat} & 2.8 & 2.9 & 1.9 \\
\bottomrule
\end{tabular}
\end{center}

CE and Dice saturate at that magnitude and the gradient overflows. Nothing clips
or normalises the product. This is a property of the method as released, not of
this hardware.

\textbf{\code{concat} rebuilds a random convolution on every forward pass.} The
$1\times1$ conv that fuses the concatenated maps is constructed \emph{inside}
\code{forward}, randomly initialised each call, never registered as a submodule
and never given to the optimizer. Calling the loss twice with identical inputs
under \code{no\_grad}:

\begin{center}
\begin{tabular}{@{}lr@{}}
\toprule
\textbf{Operation} & \textbf{Max difference between two identical calls} \\
\midrule
\code{add} / \code{mul} / \code{wf} & 0.0 (deterministic) \\
\textbf{\code{concat}} & \textbf{5.41 (non-deterministic)} \\
\bottomrule
\end{tabular}
\end{center}

None of the loss module's 66 parameter tensors belongs to it. So the
\code{concat} branch cannot learn and injects a fresh random projection into the
decoder each step. This predicts exactly what the weight readout shows:
\code{concat}'s eleven learned weights collapsed to a common 0.3755--0.3770 and
stopped moving, because there is no gradient signal for them to differentiate
on. Of the four operations, one is inert by construction and another is
numerically unbounded.

\subsection{Three further defects in the released code}

\begin{itemize}
\item \textbf{The training script cannot start.}
\code{train\_synapse\_lomix.py:10} imports \code{PVT\_CASCADE}, defined nowhere
in the repository; the only other reference is a commented-out line.
\item \textbf{Training continues in eval mode.} \code{model.train()} is called
once before the epoch loop, \code{inference()} calls \code{model.eval()} after
every epoch, and nothing switches back --- so from epoch 1 the decoder's
BatchNorm uses frozen running statistics. The same pattern appears in EMCAD and
G-CASCADE. Left intact here, since the published results were produced by this
code; exposed behind an opt-in flag for a separate experiment.
\item \textbf{Windows portability.} \code{worker\_init\_fn} is a local function,
which \code{spawn} cannot pickle, so multi-worker loading fails on Windows and
works on Linux.
\end{itemize}

\section{Limitations of this replication}

Single seed, 20 of the repository's default 300 epochs, one dataset, one
encoder, no statistical testing.
The \code{deep\_supervision}/\code{mutation} gap (0.798 vs 0.779) is well inside
what one seed can produce. What these runs \textbf{do} support independently of
budget: the inference-cost claim (measured directly), the training-cost ratios
(measured under identical conditions), and the two loss defects (verified on the
model's own logits, not dependent on training length). What they do \textbf{not}
settle is whether LoMix outperforms uniform deep supervision at full budget ---
that needs the full-schedule run this hardware cannot deliver.

\section{Where I would take it}

Two directions follow directly from the findings, and the first is my Step 2
proposal.

\textbf{Uncertainty-conditioned mixing weights.} LoMix learns one global answer
to ``which scale do I trust'', then freezes it. \S4.3 shows those weights
differentiate very slowly; \S4.2 shows enumerating all 44 combinations is what
makes training expensive. Conditioning the weights on a per-sample uncertainty
estimate addresses both: it gives the weights a richer signal than a single
global scalar, and it offers a principled basis for \emph{selecting} which
subsets are worth supervising rather than enumerating all of them --- cutting
cost on precisely the axis that dominates. Because the weights live in the loss,
the inference graph stays untouched, preserving the property that makes LoMix
deployable. Evaluated on calibration (ECE, reliability) as well as Dice, this
connects directly to the calibration-driven selection in CLEAR-MoE and the
confidence-aware objective in my LiDAR detection work.

\textbf{Bounded mixing operations.} Multiplying softmax probabilities rather
than raw logits, or normalising the product, would remove the divergence mode in
\S4.4 at no representational cost --- and registering the \code{concat} conv
properly would make a quarter of the operation set functional for the first
time. Both are small changes with a measurable prediction attached: if
\code{concat} is currently inert, fixing it should change its learned weights
away from the collapsed value.

\appendix
\section*{Appendix --- slide outline (6 slides)}
\addcontentsline{toc}{section}{Appendix}

\begin{enumerate}
\item \textbf{LoMix in one slide} --- decoder logits $\rightarrow$ subset mixing
$\rightarrow$ learned weights; loss-only, therefore zero inference overhead.
\item \textbf{Why this paper} --- CLEAR-MoE (calibration-driven selection) and
PSC (self-conditioning judged on calibration) ask the same question.
\item \textbf{Setup} --- four arms, one network; the 4\,GB deviation table; AMP
measured slower, TDR watchdog as the real constraint.
\item \textbf{What reproduced} --- efficiency table; training-cost table
($2.54\times$).
\item \textbf{What did not} --- ablation table, divergence at epoch 13, and the
two loss defects with the \code{mul} magnitude and \code{concat} determinism
evidence.
\item \textbf{Next} --- uncertainty-conditioned weights, bounded operations
$\rightarrow$ Step 2 proposal.
\end{enumerate}

\end{document}
